\section{Evaluation Framework}
\label{sec:evaluation}

Because each template computes every quantity in its gold trace, the final answer and most
milestones of a response can be checked reproducibly without a model.
We check responses deterministically wherever possible, consult an LLM judge only on what the
checks leave open (\autoref{subsec:scoring}), and validate every component against domain experts
(\autoref{subsec:validation}).

\subsection{Scoring a Response}
\label{subsec:scoring}

\paragraph{Preliminaries.}
An instance comprises a question, a gold trace, and milestones $\mathcal{M}$, the values that the
template computes, the gold trace states, and the question does not give.
Its targets are the milestones in the gold final answer, or a verdict word such as a flow regime.
A response $r$ is a sequence of steps ending in a final answer.

\paragraph{Final Answer Accuracy.}
Let $u(v)$ be one unit in the last digit that a value $v$ displays, $\mathcal{C}$ a fixed set of
unit factors (e.g., m to mm), and $\epsilon = 0.002$.
A value $\hat{y}$ in the final answer matches a numerical target $y$ if, for some
$c \in \mathcal{C}$,
\begin{equation}
    |c\,\hat{y} - y| \le \max\bigl(\epsilon\,|y|,\; c\,u(\hat{y}),\; u(y)\bigr),
    \label{eq:match}
\end{equation}
so a correct rounding at either value's precision matches.
A response scores $v(r) = 1$ if every target matches, $0.5$ if some do, and $0$ if none does or no
final answer is readable, since answering within the output limit is part of the task.
\emph{Final Answer Accuracy} (FAC) is the mean of $v(r)$ over instances.

\paragraph{Milestone Coverage.}
A milestone is reached if the response states its value within 0.5\%, the templates' display
precision rather than a fitted constant, under a unit factor in $\mathcal{C}$, anywhere and in any
order.
Unmatched milestones go to an LLM judge from outside the evaluated families, \texttt{MiMo-V2.5-Pro},
which sees the question, the response, and each milestone's name and value, not the gold trace, and
rules each one reached, not needed, or missing.
We credit only reached, as the judge also excuses unstated values as not needed
(\autoref{appendix:scoring}).
With $\mathcal{R}(r)$ the milestones matched or ruled reached,
\begin{equation}
    m(r) = |\mathcal{R}(r)| \,/\, |\mathcal{M}|,
    \label{eq:coverage}
\end{equation}
and \emph{Milestone Coverage} (MC) is the mean of $m(r)$ over instances with milestones.
The judge decides 11\% to 25\% of milestones, depending on the model.
MC measures how far a derivation gets, not whether it is free of error.

\paragraph{Step-Level Diagnostics.}
Two diagnostics look for flawed steps behind a correct answer.
The arithmetic check recomputes each displayed calculation from its printed operands and flags a
result that is not a correct rounding at its printed precision.
The judged step check sends every step the arithmetic check does not flag to the same judge, with
the question and the gold steps, and flags a step labeled a conceptual or calculation error.
\emph{Arithmetic flags} and \emph{judged step flags} are the shares of correct-answer responses with
at least one such flag, reported with their precision.
\autoref{appendix:scoring} gives the match rule in full, both judge prompts, the arithmetic rule,
and the thresholds.

\subsection{Validation Against Expert Judgment}
\label{subsec:validation}

\paragraph{Agreement with Domain Experts.}
We validate each component on 300 responses from five LLMs outside the eleven we evaluate, to 60
instances of 15 templates across all branches and levels, which three domain experts of the
instance's branch label step by step.
The final-answer check agrees with the experts' three-way verdict on 0.927 of all 300 responses, and
with their verdict on 0.982 of the 281 responses they do not mark partial, where a step-matching
design with a panel of LLM judges agrees on 0.747; with $\epsilon$ fitted on either half, the
three-way agreement is 0.908 and 0.939 on the other half.
Milestone matching reaches F1 0.923 without the judge and 0.958 with it, and over all steps, the
judged step check with the arithmetic check finds incorrect steps with precision 0.707 and recall
0.603.
On the evaluated models' responses, a domain expert confirms 171 of 189 sampled arithmetic flags
(precision 0.905), and two domain experts per item confirm the final-answer check's correct
verdicts on the five strongest models in 142 of 150 readings and its incorrect ones in 83 of 98, and
call its partial ones fully correct in 34 of 52.

\paragraph{The Limit of Verification.}
Since 175 of the 178 incorrect steps that the experts find behind correct answers are arithmetic, we
plant one defect in each of 120 clean responses: 60 one-digit slips, and 60 conceptual defects that
misstate a criterion, rule, or formula without changing a digit.
No deterministic check detects a conceptual defect (0 of 60), and four LLM judges catch 13\% to 37\%
of them without false alarms, so LLM judges catch at most about a third of misstated rules behind a
correct answer.

\paragraph{Judge Independence.}
The judge shares no family with an evaluated model, as LLM judges can favor their own
family~\citep{panickssery2024selfrecognition, li2026preferenceleakage}; re-judging 220 sampled
responses with \texttt{Grok 4.6}, from a third family, changes each model's coverage of them by
$-0.017$ to $+0.020$.
Absent self-preference is no proof of accuracy: each judge of the step-matching design's panel
favors no family, yet accepts most of the incorrect steps it judges.
\autoref{appendix:validation} details the study, the compared designs, the planted defects, and
judge independence.
